
\chapter{Post-processing}
\pgst{tfemch}

\label{ch:postprocessing}

\def\chaptertitle{Post-processing}
\def\headdate{27th April 2012}

\tfem\ has some basic post-processing facilities for extracting and
manipulating data in sysvectors and vectors:
\begin{description}
   \item[transferring:] transfer (part of) the data in a sysvector or vector to 
   another sysvector or vector.
   \item[subscripting:] generate vector subscripts pointing to a subset of the
    data in a vector in order to extract or manipulate values.
   \item[deriving:] derive vectors and sysvectors using element subroutines.
   \item[sampling:] find the data in sysvectors and vectors in points, on
    curves, surfaces or objects. 
   \item[integration:] generate a small number of integrated quantities over
    the internal domain or the boundary of a domain (curves, surfaces). For
    example: flow rate, drag, bulk stress.
\end{description}
 For
graphical post-processing we need add-ons and/or external programs.

\section{Transfer data between vectors}

A common use of transfer of data is to generate a vector containing a
single physical quantity. For example, in the driven cavity problem in
Sec.~\ref{sec:driven_cavity2D} the velocity and pressure are extracted from the
system vector using the routine \texttt{extract\_physvector}. 
\begin{listing}{1}

    call extract_physvector ( mesh, problem, sysvector, physvector )

\end{listing}
The inverse
routine, transferring a single physical quantity to a system vector, is also
available:
\begin{listing}{1}

    call store_physvector ( mesh, problem, sysvector, physvector )
 
\end{listing}

The routine \texttt{extract\_vector} has been designed to transfer data between
vectors.
It is also possible to reduce vectors by extracting only part of the data to a
new vector. For example:
\begin{listing}{1}

  call extract_vector ( mesh, problem, velocity, v_x, indegfd=[1] )

\end{listing}
will transfer the first degree of freedom of the vector \texttt{velocity} to the
first degree of freedom of a new vector \texttt{v\_x}. Note that:
\begin{itemize}
   \item both the input and output vector must already have been created,
   \item both the degrees of freedom of the input and output vector can be
specified,
   \item all degrees of freedom specified must exist, except for nodes that
         have zero degrees of freedom.
\end{itemize}
The latter property of the routine makes it easy to extract, for example, only
the values in the vertices to a new vector.

The most general routine for transferring data between sysvectors and/or vectors
is \texttt{transfer\_data}. Note, that:
\begin{itemize}
   \item transfer between different problems is possible,
   \item transfer from any sysvector/vector to any other sysvector/vector is
possible,
   \item restricting the data to physical quantities, degrees and a layer is
possible,
   \item the transfer of data is only done where possible and no data is
``invented''. 
   \item transfer is done ``node for node'' and can be somewhat slower than the
other routines above.
\end{itemize}
For example:
\begin{listing}{1}

  call transfer_data ( mesh, problem, problem2, &
    sysvector1=sol, vector2=vec2, physq1=[1], layer1=1 )

\end{listing}
will transfer layer 1 of physical quantity 1 in sysvector \texttt{sol} of
problem \texttt{problem} to
vector \texttt{vec2} of problem \texttt{problem2} (which might not have layers
defined). Note, that in nodes where layer 1 is not present, no data is
transferred and data in the vector \texttt{vec2} might not be properly defined
in some nodes. It is the responsibility of the user that all data in a vector is
properly defined by, for example, properly initializing a vector before
transferring data.

\section{Vector subscripts}
\label{sec:vectorsubscript}

Vector subscripts are one-dimensional integer arrays in \ftn\ 
that selects a section of a large array. In \tfem\ we have defined two
new types
\texttt{subscript\_t} and \texttt{subscriptcon\_t} 
that contain an
integer array component \texttt{s} that can be used as a vector subscript in
the data component \texttt{u} of a sysvector. The first type points to
nodal point degrees of freedom, whereas the second type points to
degrees of freedom related to the constraints.
A third type \texttt{subscriptvec\_t} has been defined
that contains an
integer array component \texttt{s} that can be used as a vector subscript in
the data component \texttt{u} of a vector.

In Sec.~\ref{sec:solidbar} an example is given where vector subscripts are
created for the displacement and the pressure in the solution vector. In an
increment only the displacements are updated. Another example is given in
the problem discussed in Sec.~\ref{sec:bcfcouette}:
\begin{listing}{1}

  type(subscript_t) :: vely
  type(subscriptvec_t) :: cxx, cxy, cyy
  type(vector_t) :: ctensor

  ...

! create a vector subscript for the vertical velocity for post processing

  call create_subscript ( mesh, problem, vely, physqarr=[physqvel], degfd=2 )

  ...

! subscripts and vector for ctensor

  call create_subscript ( mesh, problemq, cxx, degfd=1, vec=3 )
  call create_subscript ( mesh, problemq, cxy, degfd=2, vec=3 )
  call create_subscript ( mesh, problemq, cyy, degfd=3, vec=3 )

  call create_vector ( problemq, ctensor, vec=3 )

  ...

!  write monitor data

   write(unit=13,fmt=*) &
     step, step*deltat, maxval(sol%u(vely%s)), minval(sol%u(vely%s)), &
     sqrt( &
     sum( ( ctensor%u(cxx%s) - csteady(1) )**2 + &
          ( ctensor%u(cxy%s) - csteady(2) )**2 + &
          ( ctensor%u(cyy%s) - csteady(3) )**2 ) / 3 / size(cxx%s) )

\end{listing}
Here the velocity component in $y$-direction and the three tensor components
$c_{xx}$, $c_{xy}$ and $c_{yy}$ are selected for further processing using
vector subscripts. Note, that vector subscript act directly on the data itself.

Various heading parameters exists to limit the vector subscript to degrees of
freedom in specified points, curves, surfaces, nodesets, elementsets or groups of elements
(\texttt{points}, \texttt{curves}, \texttt{surfaces}, \texttt{nodesets}, \texttt{elementsets},
\texttt{groups})
or to exclude degrees of freedom from
certain points (\texttt{excludepoints}, \texttt{excludecurves},
\texttt{excludesurfaces}, \texttt{excludenodesets}, \texttt{excludeelementsets},
\texttt{xmin}, \texttt{xmax}).

An example using a vector subscript for getting constraint data from
the sysvector is \texttt{stokes17} (see Sec.~\ref{sec:rotatingobject}).

\section{Deriving vectors}

A general way of extracting (derived) data from sysvectors and vector is
the routine \texttt{derive\_vector}. It is used in about every example in this
userguide and uses element subroutines to generate complicated data, such as
gradients. By default it performs averaging of values at
common nodes between elements, however this can be turned off. In the latter
case (no averaging), all individual element data can be extracted by defining
the computed vector as stored `elementwise':
\begin{listing}{1}
   
  call create_vector ( problem, vector, vec=1, elementwise=.true. )

\end{listing}
The vector defined elementwise can be plotted using Figplot
(see Sec.~\ref{sec:figplot}) and sampled (see Sec.~\ref{sec:sample}), but
usually cannot be handled by external graphics packages\footnote{Except gmsh,
 see Sec.~\ref{sec:gmshwrite}.}.

A routine that is similar to \texttt{derive\_vector} but now producing a system
vector is \texttt{fill\_sysvector\_elem}. It cannot do averaging.

A `bare bones' approach of the element by element filling of data vectors is
the routine \texttt{loop\_over\_elements}. All input/output must be done by the
user via the structure \texttt{old\_vectors}. Note that there are routines
available, starting all with \texttt{get\_}, to obtain coordinates,
data and position of the data in sysvectors and vectors (see
Sec.~\ref{sec:elemsub}). In this way, the user is able to extract and store
data in the structure \texttt{old\_vectors}.

\section{Sampling}
\label{sec:sample}

Sampling is extracting data from the sysvectors and vectors in a number
of sampling points. For sampling you need to define a sample, which is of the
type \texttt{sample\_t} (see source code with \texttt{tfvi sample\_t}):
\begin{listing}{1}

  type(sample_t) :: sample

\end{listing}
The sample needs to be filled with data from nodes, points, a curve, surface or
object using the subroutine \texttt{fill\_sample}.
Additional curves, objects for
sampling can be defined using the routine \texttt{add\_to\_mesh} (see
Sec.~\ref{sec:addtomesh}).

\subsection{Sampling nodal values using points, curves, \ldots}

Geometrical entities like points and curves are linked to nodal points. Since
vectors (of type \texttt{sysvector\_t} and \texttt{vector\_t}) are defined in
nodal points, the data in these vectors can be sampled directly from the nodal
point data. For example:
\begin{listing}{1}
  call fill_sample ( mesh, problem, sample, ndegfd=1, curve=2, vector=pressure )
\end{listing}
will fill the sample using the vector \texttt{pressure} of type
\texttt{vector\_t} and
\begin{listing}{1}
  call fill_sample ( mesh, problem, sample, ndegfd=1, physq=2, &
    points=[3, 5], sysvector=sol )
\end{listing}
will sample the physical quantity 2 from solution vector \texttt{sol} in the
points \texttt{P3} and \texttt{P5} using the first degree of freedom of that
quantity. It is also possible to sample by using nodal point numbers:
\begin{listing}{1}
  call fill_sample ( mesh, problem, sample, ndegfd=2, physq=1, &
    nodes=[1336, 12999, 22], sysvector=sol )
\end{listing}
will sample the physical quantity 1 from solution vector \texttt{sol} in the
the nodal points 1336, 12999 and 22 using the first two degrees of freedom.

The number of samples in the structure sample (after sampling) is in
\texttt{sample\%nnodes}, the coordinates of the sample are in
\texttt{sample\%coor(:,:)} and the data in \texttt{sample\%u(:,:)}. The first
index of both arrays is the sample number.

\subsection{Sampling objects}
\label{sec:samplingobjects}

Objects have points that can be anywhere in the mesh. Sampling data from
vectors needs to be performed by interpolation. Since \tfem\ doesn't know
anything about the actual \fem\ shape functions being used for the data values
in the vector data, the user must supply an interpolating subroutine in a
similar way as an element subroutine in the matrix build step.

In the example file \texttt{stokes9.f90}, which is an extension of
\texttt{stokes8.f90} (use \texttt{getexample stokes}) two different samples
have been used: one single point which is sampled at each time step and points
on a line which is sampled only once.

\begin{listing}{120}
  integer, parameter :: nc = 101 ! number of nodes in cross section
  real(dp) :: xs(1,2), xc(nc,2)
                                                                                
\end{listing}
The coordinate arrays for sampling are defined here.
\begin{listing}{189}

! one object for sampling the velocity in a single point

  xs(1,:) = [0.5_dp,0.85_dp]

  call add_to_mesh ( mesh, object='coordinates', nnodes=1, coor=xs )

! one object for sampling the velocity on a line (diagonal)

  xc(:,1) = [ (i*0.01_dp, i=0,nc-1) ]
  xc(:,2) = [ (i*0.01_dp, i=0,nc-1) ]

  call add_to_mesh ( mesh, object='coordinates', nnodes=nc, coor=xc )

\end{listing}
Here two new objects are defined for sampling. Note, that there is already one
object defined earlier (the particle).
\begin{listing}{340}

!   sample in one point

    oldvectors%s(1)%p => sol

    call fill_sample ( mesh, problem, sample1, ndegfd=2, object=2, &
      elemsub=stokes_sample_velocity, coefficients=coefficients, &
      oldvectors=oldvectors )

    write( unit=10, fmt=* ) time, sample1%u(1,1), sample1%u(1,2)

\end{listing}
At line 343 the solution vector is directed towards the subroutine
\texttt{stokes\_sample\_velocity} using a pointer in oldvectors.
The sample is filled and
subsequently written to a file. This happens at each time step. So after the
program stops the value of the velocity in a single point is written to the
file as a function of time.
\begin{listing}{357}

! sample velocity on a diagonal line

  oldvectors%s(1)%p => sol

  call fill_sample ( mesh, problem, sample2, ndegfd=2, object=3, &
    elemsub=stokes_sample_velocity, coefficients=coefficients, &
      oldvectors=oldvectors )

  open( unit=10, file='sample2.out' )
  do i = 1, nc
    write( unit=10, fmt=* ) sample2%coor(i,1), sample2%coor(i,2), &
                            sample2%u(i,1), sample2%u(i,2)
  end do
  close( unit=10 )

\end{listing}
Here the sample on a line is created, filled and written to a file.
\begin{listing}{409}
  call delete ( sample1, sample2 )   
\end{listing}
The samples are released from memory.

\red{WARNING}: if the co-ordinates of an object are
changed directly by the user, the reference coordinates of the object
points\footnote{
\texttt{mesh\%objects(:)\%refcoor} and for two-sided objects
also \texttt{mesh\%objects(:)\%refcoor2}}
and the element and group information 
of the object points\footnote{\texttt{mesh\%objects(:)\%grpelm}
and for two-sided objects also \texttt{mesh\%objects(:)\%grpelm2}}
are out of sync.
In that case,
the routine \texttt{find\_refcoor\_objects} must be called for each
affected object to make the data up-to-date. Note, that
\texttt{find\_refcoor\_objects} is called in the routine
\texttt{fill\_mesh\_parts}. So that when you make a new mesh and finalize the
mesh with \texttt{fill\_mesh\_parts} the object data is up-to-date.

\red{WARNING}: the routine \texttt{find\_refcoor\_objects} uses blocks (see
Sec.~\ref{sec:blocks}) to find the elements and reference co-ordinates.
This means that if the coordinates of the mesh have been
changed directly by the user (e.g.\ when an ALE mesh is used),
the bounds of the blocks need to be updated before calling
\texttt{find\_refcoor\_objects}:
\begin{listing}{1}

  call find_bounds_blocks ( mesh )
  call find_refcoor_objects ( mesh )

\end{listing}
It is also possible to fully redefine the blocks (see Sec.~\ref{sec:blocks}),
but this takes more time.
If only the coordinates of the objects change,
there is no need to update the blocks.

\subsection{Avoiding deallocate/allocate of sample data in \texttt{fill\_sample}}

In the routine \texttt{fill\_sample} the sample memory is allocated on demand.
If the previously allocated memory of the sample is too small the memory is
deallocated and new memory is allocated. If this becomes time critical, it is
better to define a big enough sample from the start by using the 
routine \texttt{create\_sample} before the first call of \texttt{fill\_sample}.
For example
\begin{listing}{1}
  call create_sample ( mesh, sample, ndegfd=3, nnodes=5000 )
\end{listing}
which will create a sample that can accommodate 5000 sampling points with
3 degrees of freedom per point. If the actual number of sampling points and the
number of degrees of freedom needed is less than
or equal to 5000 and 3, respectively, no data is deallocated/allocated in the
call of \texttt{fill\_sample}.

\subsection{Sampling all nodes}

Samples can be invalid:
\begin{itemize}
   \item a sample point is outside the domain or outside a layer (if
         specified).
   \item the number of degrees of freedom \texttt{ndegfd} is larger than the
number of degrees available in a node.
   \item the number of degrees of freedom \texttt{ndegfd} is not equal to the
number of degrees specified with \texttt{degfd}.
\end{itemize}
By default, invalid samples are just skipped and no sample is taken. This
means, that the number of samples can be different from the number of nodes in
the object, curve, etc. and the correspondence of the samples with nodes is
lost. To avoid that, an optional keyword \texttt{allnodes} is available for
the routine
\texttt{fill\_sample} (and also \texttt{create\_sample}). By setting the
keyword to \texttt{.true.} the number of samples will be equal to the number of
nodes and the sequence will be the same as in the object, curve etc. Whether a
sample is valid is indicated by the logical array \texttt{valid} in the
structure \texttt{sample}.

\section{Integration}
\label{sec:integration}

Two types of integration are available:
\begin{itemize}
   \item Integration over the internal domain of the mesh 
         (internal elements) using the subroutine
         \texttt{integrate}. An element subroutine must be provided for
         integration on a single element.
         The routine \texttt{integrate} is similar to the
         \texttt{build\_system} for assembling the system, but now the element
         vectors are just added.
   \item Integration over the boundaries  using the subroutine
         \texttt{integrate\_boundary\_elements}.
         An element subroutine must be provided for integration on a single 
         boundary element.
         The \texttt{integrate\_boundary\_elements}
         routine is similar to the routine \texttt{add\_boundary\_elements},
         but now the element vectors are just added.
   \item Integration over an object using the subroutine
         \texttt{integrate\_object}.
         An element subroutine must be provided for contribution of 
         a single integration point of a single element of the object.
\end{itemize}
For example:
\begin{listing}{1}

    call integrate_boundary_elements ( mesh, problem, resultsum, &
      elemsub=stokes_flowrate, curve=2, coefficients=coefficients, &
      oldvectors=oldvectors )

\end{listing}
will compute the flow rate through curve 2
(output value in \texttt{resultsum(1)}) using the element subroutine
\texttt{elem\_flowrate} available in the module 
\texttt{stokes\_elements\_m} of add-on stokes
(see Sec.~\ref{sec:stokesaddon}).

\section{Graphical post-processing}

The core of \tfem\ doesn't have any graphical routines.

For 1D plotting
(xy-plots) you need an external program, such as Matlab or
the free gnuplot program. For interfacing to Matlab and
gnuplot the module printtofile (part of the add-on IO-utils) can be used 
(see Sec.~\ref{sec:printtofile}).

For 2D and simple 3D
plotting there is an add-on figplot (see Sec.~\ref{sec:figplot}).

For advanced 2D and 3D graphical post-processing an external program is needed,
At the time of writing interfaces to the following
programs are available as an add-on in \tfem:
\begin{itemize}
\item
The freely available Paraview program (see Sec.~\ref{sec:vtk}).
\item
The commercial Tecplot program (see Sec.~\ref{sec:tec}).
\item
The freely available Gmsh program (see Sec.~\ref{sec:gmshwrite}).
\end{itemize}

